\chapter{The Continuation Problem}
\label{ch:continuation-problem}

\begin{quotation}
\noindent\textit{Synopsis.} This chapter introduces the question motivating
the entire monograph: how can a future consumer continue a task when only a
finite artifact remains available? Rather than asking whether a substrate
represents its source faithfully in the abstract, the chapter shifts
attention to continuation --- the practical question of what actions remain
possible given a surviving substrate. Chess, painting, and media pipelines
are introduced as the recurring examples that generated the framework,
before the chapter states its central thesis: no finite substrate can be
guaranteed sufficient for arbitrary future continuation across an open task
space, because future requirements cannot be fully anticipated at the time
the substrate is produced.
\end{quotation}

\section{A Motivating Scenario}

An artifact is produced. A consumer arrives later, possibly much later,
possibly through an intermediary the producer never anticipated, and
must act on the basis of what is available to them. The question this
monograph asks is simple to state and difficult to resolve: what must
survive the interval between production and use for that later action to
succeed?

\section{Informal Statement of the Thesis}

\begin{principle}[Global Insufficiency, informal form]
\label{principle:global-insufficiency}
No finite substrate can be guaranteed sufficient for arbitrary future
continuation across an open task space.
\end{principle}

This statement is deliberately informal at this stage, and it is worth
saying plainly what kind of book follows from taking it seriously. It is
not a claim that information is destroyed, that communication is
impossible, or that faithful transmission never occurs -- Chapters
\ref{ch:chess} and~\ref{ch:proofs} both exhibit clean cases where a
substrate is fully adequate for its purpose. It is a narrower and more
tractable claim: that adequacy is always relative to a task, that the
relevant task cannot always be known in advance, and that this gap has a
precise structure worth studying rather than a vague one worth lamenting.
Chapter~\ref{ch:admissibility} gives the informal statement above a
precise form as \cref{prop:task-unavailability}, and
Part~\ref{part:budget-relay} extends it from a single substrate to a
chain of substrates, where the monograph's central new theorem is
proved (\cref{thm:composition-failure}).

\section{Continuation, Not Representation}

Continuation is defined operationally throughout this monograph: given
only a substrate $r$, the set of actions a consumer can correctly
perform in pursuit of a task $\tau$. This is deliberately narrower than
asking whether $r$ is a faithful, meaningful, or complete representation
of some original object $e$ in the abstract. Faithfulness-to-an-original
is not directly measurable by a consumer who never had access to $e$;
sufficiency-for-a-task is.

\section{Three Motivating Domains}

Three domains recur throughout the monograph and are introduced here
only informally, to be formalized in Part~\ref{part:worked-domains}. A
chess position illustrates a domain with closed, enumerable hidden
state: what is missing from the visible board (castling rights, en
passant eligibility, the repetition count) can in principle be listed
exhaustively and restored. A painting illustrates the opposite case: a
static configuration, such as a branch bent as if by wind, is compatible
with an open-ended and non-enumerable set of causal histories, and no
finite augmentation closes that space. A media pipeline -- audio
rendered into a transcript, a transcript compressed into a summary --
illustrates a domain in which loss accumulates across a sequence of
independently produced renderings, each locally reasonable, with no
shared record of what any earlier stage discarded.

These three domains are not illustrations selected after the fact; they
are the cases that generated the framework, and it is worth being
explicit about what that means in practice, since it bears on how the
rest of the book should be read. More than one definition in this
monograph was revised after failing to survive contact with a worked
numerical instance: the discard rule underlying the budget relay
(\cref{def:budget-discard}) turned out to admit two genuinely different
constructions only once a concrete example was pushed hard enough to
expose the difference (\cref{rem:discard-rules-diverge}), and a claimed
theorem about cumulative loss (\cref{prop:monotonicity}) turned out, once
stated against a concrete case, to be an immediate consequence of a
definition rather than an independent result requiring its original,
stronger label. The worked examples of Part~\ref{part:worked-domains}
are load-bearing for this reason: they are not decoration added to a
finished theory, they are part of how the theory was checked.

\section{An Accurate Roadmap}

Part~\ref{part:foundations} develops the vocabulary needed to state
\cref{principle:global-insufficiency} precisely: distinguishability loss
(\cref{ch:distinctions}), then a chapter situating that vocabulary within
the wider research program it descends from
(\cref{ch:four-operations}), checking prior claims of informal
inheritance against primary source text where available and marking
plainly where only weaker evidence could be had. Admissible continuation
(\cref{ch:admissibility}) and repair economics
(\cref{ch:repair-economics}) follow, before Chapter~\ref{ch:synchronic-to-diachronic}
identifies exactly where this imported machinery stops applying and
proves, rather than merely asserts, that no coordination mechanism
salvages it without genuinely new construction
(\cref{prop:no-global-conserved-quantity}). Part~\ref{part:budget-relay}
introduces the budget relay, the monograph's one new formal object, and
proves the Composition Failure Theorem
(\cref{thm:composition-failure}) against an explicit, fully worked
witness rather than only in the abstract. Part~\ref{part:modes-of-access}
resolves the preservation/reconstitution question raised by the informal
motivating conversation behind this project, correcting its initial
binary framing into a single continuous variable
(\cref{def:fidelity}), and states plainly, rather than papering over,
the one dependency this correction cannot settle on its own: whether
memory and representation stand in the relationship MEM|8 suggests
(Chapter~\ref{ch:mem8}). Part~\ref{part:worked-domains} returns to
chess, painting, and the rest with the full apparatus available, closing
with a numerically explicit pipeline witness
(\cref{prop:pipeline-witness}) built from the same machinery as the
book's central proof. Part~\ref{part:synthesis} closes with a single
statement tying distinguishability, admissibility, repair budgets, and
continuation together, and an explicit accounting of what remains
unresolved.

Six appendices support this argument without carrying its weight.
Appendix~\ref{app:proof} gives the Composition Failure Theorem's full
proof; Appendix~\ref{app:notation} collects the book's notation in one
place; Appendix~\ref{app:comparison} establishes, concretely rather than
only structurally, that this monograph's central result is independent
of \emph{Conservation of Ambiguity}'s; Appendix~\ref{app:extended-examples}
extends the framework to domains not given full chapters;
Appendix~\ref{app:simulation} sketches, without implementing, how the
framework could be tested computationally; and
Appendix~\ref{app:open-questions} consolidates every conjecture and open
question raised anywhere in the book into a single register, so that
what remains unproved is never left for a reader to discover piecemeal.

\section{A Note on How to Read the Labels}

This monograph distinguishes proved results from conjectured ones
throughout, and does so more insistently than is typical, because the
central claim of the book is partly a claim about the limits of what any
finite substrate -- including this one -- can establish once and for
all. Exactly one result in the entire book is new: the Composition
Failure Theorem this monograph exists to prove
(\cref{thm:composition-failure}). Nine further results carry the label
Theorem on the strength of being imported wholesale, with correct
attribution, from a source text that itself labels them Theorem --
including the Redistribution Theorem
(\cref{thm:redistribution}) and, at varying tiers of verification
described in the front matter and in Chapter~\ref{ch:four-operations},
seven theorems from Distinguishability Geometry, Persistence Before
Truth, and Observability Is Restorability, plus one from Admissibility
Before Truth transcribed rather than independently verified. Every other
substantial claim is a Proposition, a Corollary, a Conjecture, or an
Open Question, and the front matter's Note on Epistemic Status explains
the difference precisely. A reader who wants to know, at any point,
exactly how much weight a given claim can bear should look at its label
before its rhetoric.
